\section{Results}\label{sec:results}

All numbers are means over three seeds on the 292{,}052 test half-hours of the \NumHouseholds{} households. The
standard deviations over seeds are given in the tables. Gaps are expressed relative to the centralised test MAE.

\subsection{Accuracy of the arms}\label{sec:res-main}
Table~\ref{tab:main} compares all arms. Every LSTM beats both naive references. Persistence (\PersistenceMAE{}\,kWh) is a much stronger
baseline than seasonal naive (\SeasonalNaiveMAE{}\,kWh) at a 30-minute horizon, because consecutive half-hours are
strongly correlated. The centralised LSTM is the most
accurate arm (MAE \EtwocentralizedMAE{}\,kWh). It is better than local-only training (\EonelocalMAE{}\,kWh) for
\StEtwocentralizedvsEonelocalBetter{}\,\% of households (Holm-adjusted $p$ \StEtwocentralizedvsEonelocalP{}), which
confirms that cross-household data are valuable for this task.

Federated training recovers only part of that value. FedAvg with random grouping reaches \EthreefedavgiidMAE{}\,kWh,
\EthreefedavgiidGap{}\,\% above centralised, and it is worse than centralised for every household. Under moderate
heterogeneity ($\alpha=0.3$) FedAvg reaches \EfourfedavgmoderateMAE{}\,kWh and is worse than local-only training on
pooled MAE. Per household it beats local-only for only \StEfourfedavgmoderatevsEonelocalBetter{}\,\% of homes
($p_{\mathrm{Holm}}$ \StEfourfedavgmoderatevsEonelocalP{}). On RMSE the ordering reverses: FedAvg
(\EfourfedavgmoderateRMSE{}\,kWh) beats local-only (\EonelocalRMSE{}\,kWh). The shared model handles large deviations
better, while per-home models fit the typical half-hour better. Our protocol therefore reproduces the observation of
Briggs et al.~\cite{briggs2022} that plain FL can lose to local models on LCL, but only for MAE, the absolute-error
criterion.

\begin{table}[t]
\centering\scriptsize
\caption{Test error of all arms (kWh per half-hour; mean $\pm$ s.d. over three seeds). MAPE$^\ast$ is computed over
half-hours with load $\geq0.05$\,kWh. Traffic is the cumulative upload plus download over all rounds; for the
centralised arm it is the raw training-data volume.}\label{tab:main}
\begin{adjustbox}{max width=\textwidth}\begin{tabular}{lrrrrrr}
\toprule
Method & MAE & RMSE & MAPE$^\ast$ & WAPE & Best rd. & Traffic (MB) \\
\midrule
Persistence (naive) & 0.0890 & 0.1987 & 41.1 & 37.6 & n/a & 0.0 \\
Seasonal naive (t--48) & 0.1332 & 0.2773 & 63.6 & 56.3 & n/a & 0.0 \\
Local-only & 0.0840 $\pm$ 0.0003 & 0.1779 $\pm$ 0.0006 & 37.3 $\pm$ 0.4 & 35.5 & n/a & 0.0 \\
Centralised & 0.0811 $\pm$ 0.0003 & 0.1718 $\pm$ 0.0015 & 36.7 $\pm$ 0.8 & 34.3 & n/a & 14.0 \\
FedAvg, random & 0.0854 $\pm$ 0.0008 & 0.1758 $\pm$ 0.0006 & 39.0 $\pm$ 0.6 & 36.1 & 22.0 / 30 & 82.6 \\
FedAvg, $\alpha$=0.3 & 0.0860 $\pm$ 0.0003 & 0.1757 $\pm$ 0.0005 & 39.7 $\pm$ 0.2 & 36.4 & 24.0 / 30 & 82.6 \\
FedProx, $\alpha$=0.3 & 0.0858 $\pm$ 0.0001 & 0.1755 $\pm$ 0.0006 & 39.5 $\pm$ 0.1 & 36.3 & 26.0 / 30 & 82.6 \\
HA-FedProx, $\alpha$=0.3 & 0.0863 $\pm$ 0.0003 & 0.1765 $\pm$ 0.0005 & 39.6 $\pm$ 0.1 & 36.5 & 20.7 / 30 & 82.6 \\
FedAvg, cluster & 0.0857 $\pm$ 0.0003 & 0.1759 $\pm$ 0.0004 & 39.2 $\pm$ 0.2 & 36.2 & 21.4 / 30 & 82.6 \\
FedProx, cluster & 0.0858 $\pm$ 0.0002 & 0.1759 $\pm$ 0.0006 & 39.3 $\pm$ 0.2 & 36.3 & 23.3 / 30 & 82.6 \\
HA-FedProx, cluster & 0.0865 $\pm$ 0.0002 & 0.1772 $\pm$ 0.0008 & 39.5 $\pm$ 0.4 & 36.6 & 15.7 / 30 & 82.6 \\
FedAvg, $\alpha$=0.3, partial & 0.0850 $\pm$ 0.0005 & 0.1759 $\pm$ 0.0008 & 38.7 $\pm$ 0.5 & 36.0 & 22.4 / 30 & 37.2 \\
FedProx, $\alpha$=0.3, partial & 0.0857 $\pm$ 0.0003 & 0.1764 $\pm$ 0.0012 & 39.1 $\pm$ 0.3 & 36.2 & 20.7 / 30 & 41.3 \\
HA-FedProx, $\alpha$=0.3, partial & 0.0857 $\pm$ 0.0003 & 0.1769 $\pm$ 0.0011 & 39.1 $\pm$ 0.2 & 36.2 & 23.0 / 30 & 41.3 \\
FedAvg, $\alpha$=0.3, fp16 & 0.0857 $\pm$ 0.0004 & 0.1755 $\pm$ 0.0007 & 39.5 $\pm$ 0.2 & 36.3 & 25.2 / 30 & 41.3 \\
FedAvg, K=100 & 0.0894 $\pm$ 0.0013 & 0.1816 $\pm$ 0.0009 & 40.9 $\pm$ 1.5 & 37.8 & 29.0 / 30 & 413.2 \\
FedProx, K=100 & 0.0893 $\pm$ 0.0013 & 0.1817 $\pm$ 0.0011 & 40.7 $\pm$ 1.6 & 37.8 & 30.0 / 30 & 413.2 \\
HA-FedProx, K=100 & 0.0901 $\pm$ 0.0005 & 0.1823 $\pm$ 0.0009 & 41.2 $\pm$ 0.8 & 38.1 & 29.3 / 30 & 413.2 \\
FedAvg, K=1, random & 0.0808 $\pm$ 0.0005 & 0.1710 $\pm$ 0.0003 & 36.7 $\pm$ 0.3 & 34.2 & 16.3 / 30 & 4.1 \\
FedAvg, K=5, random & 0.0827 $\pm$ 0.0005 & 0.1721 $\pm$ 0.0006 & 38.1 $\pm$ 0.2 & 35.0 & 23.0 / 30 & 20.7 \\
FedAvg, K=50, random & 0.0869 $\pm$ 0.0006 & 0.1782 $\pm$ 0.0005 & 39.6 $\pm$ 0.5 & 36.8 & 26.7 / 30 & 206.6 \\
\bottomrule
\end{tabular}
\end{adjustbox}
\end{table}

Fig.~\ref{fig:pred} shows the first three test days of the household whose mean load is closest to the cohort
median. All learned models track the daily pattern and lag sharp, short spikes by one step, which is typical of
one-step forecasters of household load. The differences between arms are small relative to the volatility of the
series, which is why the paired, per-household tests below are needed to separate them.

\begin{figure}[t]
\centering
\includegraphics[width=\textwidth]{prediction_example.png}
\caption{Actual and predicted load for the household closest to the cohort's median load (chosen without reference
to any error), first three test days.}\label{fig:pred}
\end{figure}

\subsection{Where the federated gap comes from}\label{sec:res-decomp}
Table~\ref{tab:decomp} walks from centralised training to the hardest federated setting, changing one factor at a
time. Fig.~\ref{fig:frag} shows the fragmentation sweep.

Running the federated protocol with a single client that holds all 100 homes ($K=1$) reproduces centralised
accuracy: \EonefourfedavgKoneMAE{}\,kWh against \EtwocentralizedMAE{}\,kWh, with no significant per-household
difference (\StEonefourfedavgKonevsEtwocentralizedBetter{}\,\% of homes better, $p_{\mathrm{Holm}}$
\StEonefourfedavgKonevsEtwocentralizedP{}). Rounds, a fresh optimiser at every round and best-round selection
therefore cost nothing. The gap appears as soon as the same data are split. It grows steadily with the number of
clients, from \EonefourfedavgKfiveGap{}\,\% at $K=5$ through \EthreefedavgiidGap{}\,\% at $K=20$ and
\EonefourfedavgKfivezeroGap{}\,\% at $K=50$ to \EonethreefedavghouseholdGap{}\,\% when every household is its own
client (Table~\ref{tab:frag}, Fig.~\ref{fig:frag}). At $K=20$, FedAvg with random grouping is worse than the
single-client run for \the\numexpr100-\StEthreefedavgiidvsEonefourfedavgKoneBetter\relax\,\% of households.

Splitting homes over more clients also raises the measured heterogeneity between clients, so the two effects must
be separated. A regression of per-run FedAvg test MAE on $\log K$ and $H_{\mathrm{JS}}$ over all full-participation
runs ($n=\RegN{}$, $R^2=\RegRtwo{}$) attributes the error to fragmentation: each unit of $\log K$ adds
$\RegLogK{}\times10^{-3}$\,kWh ($p$ \RegLogKP{}), while the heterogeneity term is not significant ($p$ \RegHJSP{}).
Within $K=20$, where heterogeneity was dialled independently of $K$, error and $H_{\mathrm{JS}}$ are unrelated
(Spearman \KtwentySpearman{}, $p$ \KtwentySpearmanP{}, $n=\KtwentyN{}$). A direct comparison makes the same point:
50 clients of two randomly grouped homes ($H=\EonefourfedavgKfivezeroH{}$) are worse than 20 clients grouped by load
profile ($H=\EsixfedavgstrongH{}$), \EonefourfedavgKfivezeroMAE{} against \EsixfedavgstrongMAE{}\,kWh.

We attribute the effect to the averaging step. With one local epoch per round, each client takes about
$n_k/256$ optimiser steps on its own data before its model is averaged with the others. The smaller the client, the
fewer and noisier its steps, and the further the average of $K$ such models falls from the model that a single pass
over the pooled data would produce. The total amount of data and computation is the same for every $K$; what
changes is how much of it is combined by averaging weights rather than by following gradients.

\begin{table}[t]
\centering\small
\caption{Fragmentation sweep: the same 100 households split over $K$ clients with random grouping. $K=100$ is
household mode. FT: one epoch of local fine-tuning of the final global model.}\label{tab:frag}
\begin{adjustbox}{max width=\textwidth}\begin{tabular}{rrrrrr}
\toprule
$K$ & homes/client & Test MAE & FT MAE & $\bar\rho$ & Gap to central (\%) \\
\midrule
20 & 5 & 0.0853 $\pm$ 0.0009 & 0.0834 & 0.856 & +5.2 \\
\bottomrule
\end{tabular}
\end{adjustbox}
\end{table}

\begin{table}[t]
\centering\small
\caption{Decomposition of the federated--centralised gap (FedAvg, test MAE in kWh). Each row adds one factor to the
row above. Rows without a $\Delta$ are alternatives to the row above, not additions.}\label{tab:decomp}
\begin{adjustbox}{max width=\textwidth}\begin{tabular}{lrrr}
\toprule
Setting & Test MAE & $\Delta$ vs previous & Gap to central (\%) \\
\midrule
Centralised (pooled raw data) & 0.0811 $\pm$ 0.0003 & -- & +0.0 \\
+ round protocol, 1 client (K=1) & 0.0808 $\pm$ 0.0005 & -0.0002 & -0.3 \\
+ split over 20 clients, random grouping & 0.0853 $\pm$ 0.0009 & +0.0044 & +5.2 \\
+ label-skewed grouping ($\alpha$=0.3) & 0.0861 $\pm$ 0.0004 & +0.0008 & +6.2 \\
+ 50\% participation, 20\% stragglers & 0.0849 $\pm$ 0.0004 & -0.0012 & +4.7 \\
cluster grouping instead (strong non-IID) & 0.0857 $\pm$ 0.0003 & -- & +5.7 \\
1 home = 1 client (K=100) & 0.0894 $\pm$ 0.0013 & -- & +10.3 \\
\bottomrule
\end{tabular}
\end{adjustbox}
\end{table}

\begin{figure}[t]
\centering
\includegraphics[width=0.85\textwidth]{fragmentation.png}
\caption{Test MAE when the same 100 households are split over $K$ clients with random grouping ($K=100$ is one
household per client). The dashed line adds one epoch of local fine-tuning to the final global model.}\label{fig:frag}
\end{figure}

\subsection{Heterogeneity changes the price of accuracy, not accuracy}\label{sec:res-het}
Table~\ref{tab:het} and Fig.~\ref{fig:het} report the heterogeneity sweep. The load-shape divergence between clients,
$H_{\mathrm{JS}}$, increases monotonically from random grouping to the cluster partition, by a factor of about two.
The composite index $H$ does not, because random grouping produces a larger spread of client mean load than the
Dirichlet groupings, and the cluster partition, which sorts homes by load level, produces the largest spread.
Nominal labels such as $\alpha$ are therefore an unreliable proxy for how different clients actually are, which is
why we report measured values.

Across this whole range, the final FedAvg error moves by at most 0.0012\,kWh (from \EonezerofedavgaonepzeroMAE{}
to \EfourfedavgmoderateMAE{}\,kWh), comparable to the seed-to-seed standard deviation, and does not trend with
$H_{\mathrm{JS}}$ (Fig.~\ref{fig:het}a). Per household, FedAvg under the cluster partition is better than under random
grouping for \StEsixfedavgstrongvsEthreefedavgiidBetter{}\,\% of homes, and the difference is not significant after
correction ($p_{\mathrm{Holm}}$ \StEsixfedavgstrongvsEthreefedavgiidP{}). Heterogeneity does, however, change how fast
the federation gets there. The traffic needed to come within 10\,\% of centralised validation MAE grows from
\EthreefedavgiidBeps{}\,MB (\EthreefedavgiidReps{} rounds) with random grouping to \EfourfedavgmoderateBeps{}\,MB at
$\alpha=0.3$ and \EsixfedavgstrongBeps{}\,MB under the cluster partition (Fig.~\ref{fig:het}b), an increase of about
85\,\%. In this setting, heterogeneity is a communication cost rather than an accuracy cost.

\begin{figure}[t]
\centering
\includegraphics[width=\textwidth]{heterogeneity_response.png}
\caption{Heterogeneity sweep (community mode, $K=20$). (a) Final test MAE against the measured load-shape divergence
between clients; the dotted line is centralised training. (b) Megabytes exchanged before validation MAE first comes
within 10\,\% of centralised (B@$\varepsilon$). Labels give the grouping strategy.}\label{fig:het}
\end{figure}

\begin{table}[t]
\centering\small
\caption{Heterogeneity sweep: measured indices, mean update disagreement $\bar\rho$, test MAE and gap to
centralised, for each grouping and algorithm.}\label{tab:het}
\begin{adjustbox}{max width=\textwidth}\begin{tabular}{llrrrrr}
\toprule
Partition & Method & $H$ & $H_{JS}$ & $\bar\rho$ & Test MAE & Gap to central (\%) \\
\midrule
$\alpha$=0.1 & FedAvg & 0.242 & 0.0181 & 0.876 & 0.0856 $\pm$ 0.0005 & +5.6 \\
$\alpha$=0.1 & FedProx & 0.242 & 0.0181 & 0.863 & 0.0856 $\pm$ 0.0006 & +5.6 \\
$\alpha$=0.1 & HA-FedProx & 0.242 & 0.0181 & 0.844 & 0.0859 $\pm$ 0.0003 & +5.9 \\
$\alpha$=0.3 & FedAvg & 0.243 & 0.0153 & 0.874 & 0.0861 $\pm$ 0.0004 & +6.2 \\
$\alpha$=0.3 & FedProx & 0.243 & 0.0153 & 0.861 & 0.0858 $\pm$ 0.0001 & +5.9 \\
$\alpha$=0.3 & HA-FedProx & 0.243 & 0.0153 & 0.844 & 0.0863 $\pm$ 0.0003 & +6.5 \\
random & FedAvg & 0.255 & 0.0105 & 0.856 & 0.0853 $\pm$ 0.0009 & +5.2 \\
random & FedProx & 0.255 & 0.0105 & 0.845 & 0.0855 $\pm$ 0.0003 & +5.4 \\
random & HA-FedProx & 0.255 & 0.0105 & 0.822 & 0.0859 $\pm$ 0.0009 & +6.0 \\
$\alpha$=1.0 & FedAvg & 0.270 & 0.0124 & 0.867 & 0.0849 $\pm$ 0.0006 & +4.7 \\
$\alpha$=1.0 & FedProx & 0.270 & 0.0124 & 0.856 & 0.0857 $\pm$ 0.0001 & +5.7 \\
$\alpha$=1.0 & HA-FedProx & 0.270 & 0.0124 & 0.841 & 0.0857 $\pm$ 0.0008 & +5.8 \\
cluster & FedAvg & 0.381 & 0.0202 & 0.876 & 0.0857 $\pm$ 0.0003 & +5.7 \\
cluster & FedProx & 0.381 & 0.0202 & 0.866 & 0.0858 $\pm$ 0.0002 & +5.8 \\
cluster & HA-FedProx & 0.381 & 0.0202 & 0.851 & 0.0865 $\pm$ 0.0002 & +6.7 \\
\bottomrule
\end{tabular}
\end{adjustbox}
\end{table}

\subsection{Proximal regularisation and HA-FedProx}\label{sec:res-prox}
Validation-based tuning selected the smallest proximal coefficient in the grid for both methods
($\mu=\TunedMuFedProx{}$, $\mu_0=\TunedMuHA{}$). Larger values consistently increased validation error, so the data
asked for as little proximal pull as possible. With the tuned coefficient, FedProx and FedAvg are practically
indistinguishable (Table~\ref{tab:stats}). Under $\alpha=0.3$ FedProx is better for \StEfivefedproxmoderatevsEfourfedavgmoderateBetter{}\,\%
of households, but the median paired difference is \StEfivefedproxmoderatevsEfourfedavgmoderateMed{}\,kWh, about
0.1\,\% of the error. Under the cluster partition and under partial participation, FedProx is slightly worse.

HA-FedProx never improves on FedAvg. Under the cluster partition it is worse for
\the\numexpr100-\StEonezerohafedproxstrongvsEsixfedavgstrongBetter\relax\,\% of households (median
\StEonezerohafedproxstrongvsEsixfedavgstrongMed{}\,kWh, $p_{\mathrm{Holm}}$ \StEonezerohafedproxstrongvsEsixfedavgstrongP{}),
and under mild heterogeneity it needs more traffic to reach the $\varepsilon$ band
(\EonezerohafedproxiidBeps{}\,MB with random grouping against \EthreefedavgiidBeps{}\,MB for FedAvg; Fig.~\ref{fig:het}b).
Fig.~\ref{fig:rho} shows why. The update-disagreement ratio $\rho_t$ starts low (\RhoRoundOne{} on average in round 1),
rises quickly and settles near \RhoLast{} as the global model converges and the common descent direction
$\bar\Delta$ shrinks relative to the client-specific parts of the updates. The signal does rank heterogeneity
correctly: across the full-participation runs with $K=20$, mean $\rho$ increases with $H_{\mathrm{JS}}$ (Spearman
\RhoHSpearman{}, $p$ \RhoHSpearmanP{}, $n=\RhoHRuns{}$), and the round-1 value even more so (Spearman
\RhoOneHSpearman{}). However, the spread of mean $\rho$ across groupings is only \RhoBetween{}, whereas $\rho_t$
rises by \RhoWithin{} within a typical run, about \RhoWithinOverBetween{} times more. A rule
$\mu_{t+1}\propto\rho_t$ is therefore, in effect, a schedule that increases the proximal pull as training converges,
whatever the heterogeneity. Late in training this suppresses useful local progress, which is the opposite of what
the validation-tuned FedProx results ask for.

\begin{figure}[t]
\centering
\includegraphics[width=0.85\textwidth]{rho_vs_round.png}
\caption{Update-disagreement ratio $\rho_t$ per round for FedAvg under different groupings (mean $\pm$ s.d. over
seeds). The within-run rise dominates the differences between groupings.}\label{fig:rho}
\end{figure}

\begin{table}[t]
\centering\scriptsize
\caption{Paired Wilcoxon signed-rank tests over the 100 households (per-household test MAE averaged over seeds).
``A better'' is the share of households for which A has the lower MAE. $p_{\mathrm{Holm}}$ is Holm-adjusted across
all rows.}\label{tab:stats}
\begin{adjustbox}{max width=\textwidth}\begin{tabular}{llrrrr}
\toprule
A & B & median $\Delta$MAE (A$-$B) & A better (\%) & $p$ & $p_{\mathrm{Holm}}$ \\
\midrule
FedAvg, $\alpha$=0.3 & Local-only & +0.0013 & 34 & 0.0002 & 0.0023 \\
FedAvg, $\alpha$=0.3 & Centralised & +0.0026 & 0 & $<10^{-4}$ & $<10^{-4}$ \\
FedAvg, random & Centralised & +0.0021 & 0 & $<10^{-4}$ & $<10^{-4}$ \\
FedProx, $\alpha$=0.3 & FedAvg, $\alpha$=0.3 & -0.0001 & 75 & $<10^{-4}$ & $<10^{-4}$ \\
FedProx, cluster & FedAvg, cluster & +0.0001 & 36 & 0.0011 & 0.0103 \\
FedProx, $\alpha$=0.3, partial & FedAvg, $\alpha$=0.3, partial & +0.0003 & 24 & $<10^{-4}$ & $<10^{-4}$ \\
HA-FedProx, $\alpha$=0.3 & FedAvg, $\alpha$=0.3 & +0.0000 & 44 & 0.2480 & 0.4960 \\
HA-FedProx, $\alpha$=0.3 & FedProx, $\alpha$=0.3 & +0.0001 & 43 & 0.0031 & 0.0217 \\
HA-FedProx, cluster & FedAvg, cluster & +0.0004 & 27 & $<10^{-4}$ & $<10^{-4}$ \\
HA-FedProx, cluster & FedProx, cluster & +0.0001 & 39 & 0.0012 & 0.0103 \\
HA-FedProx, $\alpha$=0.3, partial & FedAvg, $\alpha$=0.3, partial & +0.0003 & 21 & $<10^{-4}$ & $<10^{-4}$ \\
HA-FedProx, $\alpha$=0.3, partial & FedProx, $\alpha$=0.3, partial & +0.0002 & 35 & 0.0469 & 0.1875 \\
FedAvg, cluster & FedAvg, random & +0.0001 & 45 & 0.0177 & 0.0884 \\
FedAvg, K=100 & Local-only & +0.0027 & 18 & $<10^{-4}$ & $<10^{-4}$ \\
FedAvg, K=100 & Centralised & +0.0045 & 0 & $<10^{-4}$ & $<10^{-4}$ \\
HA-FedProx, K=100 & FedAvg, K=100 & +0.0003 & 22 & $<10^{-4}$ & $<10^{-4}$ \\
FedAvg, $\alpha$=0.3, fp16 & FedAvg, $\alpha$=0.3 & -0.0002 & 92 & $<10^{-4}$ & $<10^{-4}$ \\
Centralised & Local-only & -0.0016 & 72 & $<10^{-4}$ & $<10^{-4}$ \\
FedAvg, K=1, random & Centralised & -0.0000 & 52 & 0.1046 & 0.3138 \\
FedAvg, random & FedAvg, K=1, random & +0.0023 & 1 & $<10^{-4}$ & $<10^{-4}$ \\
FedAvg, $\alpha$=0.3 + local FT & Local-only & -0.0004 & 52 & 0.3900 & 0.4960 \\
FedAvg, $\alpha$=0.3 + local FT & Centralised & +0.0010 & 27 & $<10^{-4}$ & $<10^{-4}$ \\
FedAvg, K=100 + local FT & Local-only & +0.0006 & 39 & 0.0142 & 0.0854 \\
\bottomrule
\end{tabular}
\end{adjustbox}
\end{table}

\subsection{Communication}\label{sec:res-comm}
One model message is \PayloadKB{}\,kB. Thirty rounds of full participation with $K=20$ clients therefore move
\EthreefedavgiidMB{}\,MB, about six times the \EtwocentralizedMB{}\,MB of raw training readings that centralised
training would upload once. The bytes needed for near-centralised accuracy are much smaller, and comparable to the
data volume: \EthreefedavgiidBeps{}\,MB with random grouping and \EsixfedavgstrongBeps{}\,MB under strong
heterogeneity (Fig.~\ref{fig:bytes}). Partial participation halves the traffic per round. FedAvg with 50\,\%
participation and 20\,\% stragglers reaches the $\varepsilon$ band after \EeightfedavgpartialBeps{}\,MB and ends at
\EeightfedavgpartialMAE{}\,kWh, no worse than full participation, using \EeightfedavgpartialMB{}\,MB in total.
 Sending updates as 16-bit floats halves the traffic (\EonetwofedavgfponesixmoderateMB{}\,MB in total,
\EonetwofedavgfponesixmoderateBeps{}\,MB to the $\varepsilon$ band) and does not harm accuracy:
\EonetwofedavgfponesixmoderateMAE{}\,kWh against \EfourfedavgmoderateMAE{}\,kWh at 32 bits. The 16-bit run is even
marginally better for \StEonetwofedavgfponesixmoderatevsEfourfedavgmoderateBetter{}\,\% of households (median
\StEonetwofedavgfponesixmoderatevsEfourfedavgmoderateMed{}\,kWh), possibly because rounding acts as a mild regulariser.

Fragmentation dominates the bill. One client moves only \EonefourfedavgKoneMB{}\,MB over 30 rounds. With
household clients the federation moves \EonethreefedavghouseholdMB{}\,MB, about 30 times the raw training data.
Only two of the three household-mode FedAvg runs reached the $\varepsilon$ band within 30 rounds; those took
\EonethreefedavghouseholdReps{} rounds and \EonethreefedavghouseholdBeps{}\,MB on average. None of the HA-FedProx
household runs reached it.

Validation error does not decrease monotonically with traffic. All federated variants pass through a transient rise
in rounds 5--8 before converging (Fig.~\ref{fig:bytes}). This transient coincides with the dip in $\rho_t$
(Fig.~\ref{fig:rho}), when client updates briefly align. It means that the round budget, and hence the traffic, must
be chosen on validation error rather than fixed in advance.

\begin{figure}[t]
\centering
\includegraphics[width=0.85\textwidth]{mae_vs_cumulative_bytes.png}
\caption{Validation MAE of the global model against cumulative traffic (upload plus download, log scale). The dotted
line is the centralised validation MAE.}\label{fig:bytes}
\end{figure}

\subsection{Who gains from federation?}\label{sec:res-equity}
Table~\ref{tab:equity} and Fig.~\ref{fig:equity} report the per-household view. Centralised training leaves
\EqEtwocentralizedHarmed{}\,\% of households worse off than training alone. FedAvg leaves \EqEthreefedavgiidHarmed{}\,\%
(random grouping) to \EqEfourfedavgmoderateHarmed{}\,\% ($\alpha=0.3$) worse off. A single epoch of local fine-tuning
of the final global model, which needs no communication, reduces this to \EqEthreefedavgiidFTHarmed{}--\EqEfourfedavgmoderateFTHarmed{}\,\%.
After fine-tuning, FedAvg is statistically indistinguishable from local training per household
(\StEfourfedavgmoderateFTvsEonelocalBetter{}\,\% better, $p_{\mathrm{Holm}}$ \StEfourfedavgmoderateFTvsEonelocalP{})
and improves pooled MAE to \EfourfedavgmoderateFT{}\,kWh. It remains behind centralised training
(\StEfourfedavgmoderateFTvsEtwocentralizedBetter{}\,\% of households better).
 With household clients the picture is worse. FedAvg is worse than local training for
\EqEonethreefedavghouseholdHarmed{}\,\% of households ($p_{\mathrm{Holm}}$
\StEonethreefedavghouseholdvsEonelocalP{}), and its 90th-percentile household error rises to
\EonethreefedavghouseholdPninetyMAE{}\,kWh, against \EonelocalPninetyMAE{}\,kWh for local training. Fine-tuning
reduces the share harmed to \EqEonethreefedavghouseholdFTHarmed{}\,\% and the pooled MAE to
\EonethreefedavghouseholdFT{}\,kWh, which is still not better than local training (\StEonethreefedavghouseholdFTvsEonelocalBetter{}\,\%
of households better, $p_{\mathrm{Holm}}$ \StEonethreefedavghouseholdFTvsEonelocalP{}).

By consumption quintile (Fig.~\ref{fig:equity}b), the advantage of pooling the raw data is concentrated in the
highest-consuming homes, whose errors are largest in absolute terms. FedAvg does not reproduce this advantage: in the
top quintile its median household error is above that of local training. The households that federation fails most
are therefore the ones whose forecasts matter most for network operation.

\begin{figure}[t]
\centering
\includegraphics[width=0.85\textwidth]{federation_benefit_cdf.png}\\[4pt]
\includegraphics[width=0.85\textwidth]{equity_by_stratum.png}
\caption{(a) Distribution over households of the benefit of each collaborative arm over local training
($\mathrm{MAE}_{\mathrm{local}}-\mathrm{MAE}_{\mathrm{arm}}$; left of zero means worse off). (b) Median household test
MAE by quintile of training-period mean load.}\label{fig:equity}
\end{figure}

\begin{table}[t]
\centering\small
\caption{Equity of federation relative to local-only training. P90 MAE is the 90th percentile of household test
MAE. FT: one epoch of local fine-tuning of the final global model.}\label{tab:equity}
\begin{adjustbox}{max width=\textwidth}\begin{tabular}{lrrrr}
\toprule
Arm & Harmed vs local (\%) & Median benefit (kWh) & P90 MAE & Harmed after FT (\%) \\
\midrule
Centralised LSTM (pooled raw data) & 28 & +0.0016 & 0.1364 & -- \\
FedAvg, random (IID-ish) & 56 & -0.0005 & 0.1439 & 46 \\
FedAvg, $\alpha$=0.3 & 66 & -0.0013 & 0.1445 & 48 \\
\bottomrule
\end{tabular}
\end{adjustbox}
\end{table}
